\documentclass[12pt,a4paper]{article}
\usepackage[left=15mm,right=15mm,top=14mm,bottom=19mm,footskip=9mm]{geometry}
\usepackage{iftex}
\ifPDFTeX
  \usepackage[T5]{fontenc}
  \usepackage[utf8]{inputenc}
  \usepackage{mathptmx}
\else
  \usepackage{fontspec}
  \IfFontExistsTF{Times New Roman}{\setmainfont{Times New Roman}}{%
    \IfFontExistsTF{Tinos}{\setmainfont{Tinos}}{\setmainfont{TeX Gyre Termes}}}
\fi
\usepackage{amsmath,amssymb}
\usepackage{tikz}
\usetikzlibrary{arrows.meta,calc}
\usepackage[most]{tcolorbox}
\usepackage{fancyhdr,lastpage}
\usepackage{tabularx,array}
\definecolor{sectionpurple}{RGB}{155,42,143}
\definecolor{notebackground}{RGB}{229,240,245}
\definecolor{noteborder}{RGB}{255,55,55}
\definecolor{graphred}{RGB}{188,55,61}
\definecolor{graphblue}{RGB}{45,95,200}
\setlength{\parindent}{0pt}
\setlength{\parskip}{1.5pt}
\setlength{\headheight}{15pt}
\setlength{\abovedisplayskip}{4pt}
\setlength{\belowdisplayskip}{4pt}
\setlength{\abovedisplayshortskip}{2pt}
\setlength{\belowdisplayshortskip}{2pt}
\pagestyle{fancy}
\fancyhf{}
\renewcommand{\headrulewidth}{0pt}
\renewcommand{\footrulewidth}{0.4pt}
\fancyfoot[L]{\fontsize{8.5}{10}\selectfont Tài liệu tổng hợp theo TDMXX22 (Tư Duy Mở)}
\fancyfoot[R]{\fontsize{8.5}{10}\selectfont Trang \thepage/\pageref{LastPage} -- Bài giảng 08.1}
\newtcolorbox{notebox}{colback=notebackground,colframe=noteborder,
  boxrule=0.8pt,arc=2pt,left=7pt,right=7pt,top=3pt,bottom=3pt,
  before skip=2pt,after skip=2pt}
\newcommand{\partheading}[1]{\par\vspace{4pt}\noindent\textbf{\textcolor{sectionpurple}{#1}}\par\vspace{1pt}}
\newcommand{\example}[2][TDM22]{\par\vspace{1pt}\noindent\textbf{\textcolor{sectionpurple}{Ví dụ #2:} \textcolor{noteborder}{[#1]}}\enspace}
\newcommand{\tightitem}{\setlength{\itemsep}{0pt}\setlength{\topsep}{1pt}\setlength{\parsep}{0pt}\setlength{\partopsep}{0pt}}
\newcommand{\veczero}{\overrightarrow{0}}
\newcommand{\cvec}[1]{\overrightarrow{#1}}
\newcommand{\cross}[2]{\left[\vec{#1},\vec{#2}\right]}
\begin{document}
\fontsize{12}{14.6}\selectfont
\begin{center}
  \textbf{BÀI GIẢNG 08.1}\\[-1pt]
  \textbf{TÍCH CÓ HƯỚNG CỦA HAI VECTƠ TRONG KHÔNG GIAN OXYZ}\\[-2pt]
  {\fontsize{10.5}{12}\selectfont\itshape Biên soạn theo chủ đề TDMXX22 -- Tư Duy Mở}
\end{center}
\vspace{-7pt}

\partheading{PHẦN I. ĐỊNH NGHĨA VÀ TÍNH CHẤT}
\textbf{1. Định nghĩa}
Trong không gian $Oxyz$, cho hai vectơ $\vec u=(a;b;c)$ và $\vec v=(a';b';c')$.
Tích có hướng (hay còn gọi là tích vectơ) của $\vec u$ và $\vec v$ là một vectơ,
kí hiệu $\cross{u}{v}$, có tọa độ là:
\[
\cross{u}{v}
=\left(
\begin{vmatrix}b&c\\b'&c'\end{vmatrix};
\begin{vmatrix}c&a\\c'&a'\end{vmatrix};
\begin{vmatrix}a&b\\a'&b'\end{vmatrix}
\right)
=(bc'-b'c;\ a'c-ac';\ ab'-a'b).
\]
\textit{Mẹo ghi nhớ:} “báo cáo, các anh, ăn bánh”.

\begin{notebox}
\example[TDM21]{01} Trong không gian $Oxyz$, cho hai vectơ
$\vec u=(1;2;3)$ và $\vec v=(4;5;6)$. Hãy tính:
\[
\cross{u}{v},\qquad \cross{v}{u},\qquad
\cross{u}{v}\cdot\vec u,\qquad
\cross{u}{v}\cdot\vec v\,?
\]
\end{notebox}

\textbf{2. Tính chất và ứng dụng cơ bản}
Trong không gian $Oxyz$, cho hai vectơ $\vec u$, $\vec v$, ta có:
\begin{itemize}\tightitem
 \item $\cross{v}{u}=-\cross{u}{v}$.
 \item $\cross{u}{v}=\veczero\iff\vec u,\vec v$ cùng phương.
 \item Nếu $\cross{u}{v}\ne\veczero$ thì $\cross{u}{v}\perp\vec u$ và $\cross{u}{v}\perp\vec v$.
\end{itemize}
\begin{notebox}
\example[TDM31]{01} Trong không gian $Oxyz$, cho ba điểm
$A(2;1;0)$, $B(4;1;2)$, $C(1;2;0)$.
\begin{enumerate}\tightitem
 \item Hãy xác định một vectơ pháp tuyến của mặt phẳng $(ABC)$.
 \item Hãy tìm tọa độ điểm $D=(ABC)\cap Oy$.
\end{enumerate}
\end{notebox}

\partheading{PHẦN II. CÁC BÀI TOÁN ỨNG DỤNG}
\textbf{3. Bài toán xác định góc trong không gian Oxyz}
\begin{notebox}
\example[TDM31]{01} Trong không gian $Oxyz$, cho các điểm
$A(2;0;0)$, $B(0;1;1)$, $C(1;0;2)$, $D(-2;0;1)$.
Hãy xác định:
\begin{enumerate}\tightitem
 \item $\cos\varphi_1$, với $\varphi_1$ là góc giữa hai mặt phẳng $(ABC)$ và $(DBC)$.
 \item $\cos\varphi_2$, với $\varphi_2$ là góc giữa đường thẳng $AD$ và mặt phẳng $(ABC)$.
\end{enumerate}
\end{notebox}

\textbf{4. Ứng dụng tích có hướng để tính diện tích và thể tích}
Trong không gian $Oxyz$, cho bốn điểm $A,B,C,D$ tạo thành một tứ diện.
Hãy ứng dụng tích có hướng của hai vectơ để tính diện tích $S_{\triangle ABC}$
của tam giác $ABC$ và thể tích $V_{ABCD}$ của tứ diện $ABCD$.
\begin{notebox}
\example[TDM31]{01} Trong không gian $Oxyz$, cho các điểm
$A(2;1;0)$, $B(4;1;2)$, $C(1;2;0)$, $D(0;2;0)$.
Hãy xác định:
\begin{enumerate}\tightitem
 \item Diện tích $S_{\triangle ABC}$.
 \item Thể tích $V_{ABCD}$.
\end{enumerate}
\end{notebox}

\textbf{5. Bài toán xác định khoảng cách trong không gian Oxyz}
Trong không gian $Oxyz$, cho tứ diện $ABCD$. Hãy xác định:
\begin{enumerate}\tightitem
 \item Khoảng cách $d(A,BC)$ từ điểm $A$ đến đường thẳng $BC$.
 \item Khoảng cách $d\bigl(A,(BCD)\bigr)$ từ điểm $A$ đến mặt phẳng $(BCD)$.
 \item Khoảng cách $d(AB,CD)$ giữa hai đường thẳng $AB$ và $CD$.
\end{enumerate}
\begin{center}
\textbf{\textcolor{sectionpurple}{--- Hết ---}}
\end{center}
\end{document}
